\chapter{State of the Art}

This chapter reviews the background needed to frame the work. It first introduces graphene and
\acrshort{gfet} biosensors for nucleic-acid detection, then surveys the biophysical principles
and computational tools used to design oligonucleotide probes, and finally the recent advances
in three-dimensional structure prediction that make an \emph{in silico} structural validation
of the probes feasible.

\section{Graphene and GFET biosensors}

Graphene is a two-dimensional allotrope of carbon whose charge carriers behave as massless
Dirac fermions, giving it very high carrier mobility and a conductance that is extremely
sensitive to the local electrostatic environment \citep{novoselov2004}. Because every atom of
graphene lies at the surface, the adsorption of even a single gas molecule was shown to be
detectable as a change in its electrical resistance \citep{schedin2007}, which established
graphene as an outstanding transducer material for sensing.

In a \acrfull{gfet}, a graphene sheet forms the conductive channel between a source and a drain
electrode, and its conductance is modulated by a gate potential. When biomolecules bind to the
functionalised graphene surface, their intrinsic charge acts as an additional local gate,
shifting the transfer characteristic of the device. For biosensing, the surface is decorated
with a bioreceptor that confers specificity. Using an \acrfull{ssdna} capture probe as the
bioreceptor turns the device into a nucleic-acid sensor: hybridisation of a complementary
target strand brings additional negative charge (the phosphate backbone) close to the channel
and produces a measurable, label-free signal. Devices of this kind have achieved ultra-low
limits of detection for nucleic acids, in some engineered geometries approaching the
single-molecule level \citep{hwang2020}, and the design, operation and performance trade-offs
of \acrshort{gfet} bioanalytical sensors have been reviewed in detail \citep{beraud2021}.

The specificity and sensitivity of the assay are ultimately governed by the capture probe,
which is the object of the computational design effort described in this dissertation.

\section{Design of oligonucleotide capture probes}

The requirements for a good hybridisation probe have been understood since the foundational
work on nucleic-acid hybridisation \citep{wetmur1991}. Three families of criteria are
generally considered.

\subsection{Hybridisation thermodynamics}

The stability of the probe--target duplex is quantified by its melting temperature
(\acrshort{tm}), the temperature at which half of the duplexes are dissociated. \acrshort{tm}
is predicted accurately with the nearest-neighbour model, whose thermodynamic parameters were
unified by \citet{santalucia2004}. A probe must have a \acrshort{tm} comfortably above the
assay temperature so that the duplex is stable during measurement, but not so high that it
loses specificity. The GC content is a related constraint: it influences both the
\acrshort{tm} and the propensity to form stable secondary structure, and a window of roughly
40--60\% is usually recommended. Software such as \texttt{primer3} \citep{untergasser2012} and
the IDT OligoAnalyzer tool \citep{idtoligo} implement these calculations and are widely used
for primer and probe design.

\subsection{Self-structure}

A probe that folds into a stable hairpin, or that dimerises with a copy of itself, is
sequestered away from its target and loses capture efficiency. The relevant quantity is the
free energy of these structures: the more negative the free energy, the more stable, and hence
more problematic, the fold. Free energies of hairpins and self-dimers are again available from
\texttt{primer3} and OligoAnalyzer, while the overall secondary-structure stability of a single
strand can be estimated from its \acrfull{mfe} using dynamic-programming folding algorithms
parameterised with nearest-neighbour energies \citep{turner2010,zadeh2011}. Lightweight
implementations such as \texttt{seqfold} \citep{seqfold} make this practical to run over
thousands of candidates.

\subsection{Sequence conservation}

To detect a pathogen across its clinical strains, a probe must target a region that is
conserved in the population. Conserved regions are identified from a \acrfull{msa} of many
sequences of the gene; MAFFT is a fast and accurate aligner widely used for this purpose
\citep{katoh2013}. Conservation can then be scored column by column. A simple measure is the
frequency of the most common base, but a more robust measure is the \acrfull{ppi}, the average
identity over all pairs of sequences in a column, which correctly handles ambiguous
(IUPAC) bases. The \acrshort{ppi} and an associated ``No-fold'' score were introduced in the
\emph{ViruScope} probe-design tool \citep{viruscope}, and are adopted in this work.

\section{Three-dimensional structure prediction}

Recent deep-learning methods have transformed the prediction of biomolecular structure.
AlphaFold demonstrated near-experimental accuracy for protein structures \citep{jumper2021},
and AlphaFold~3 extended the approach to complexes that include nucleic acids and other
molecules, together with per-residue and pairwise confidence estimates such as the
\acrfull{plddt} and the \acrfull{pae} \citep{abramson2024}. Open models such as Boltz-1 and its
successor Boltz-2 make this class of prediction freely available and runnable on commodity
GPUs \citep{wohlwend2024}. Although these confidence metrics were calibrated primarily for
proteins, the predicted structure and the \acrshort{plddt} of a short \acrshort{ssdna} probe
still provide a useful, independent check on whether a computationally selected sequence is
structurally well behaved --- an opportunity that this dissertation exploits as a final
validation stage.

\section{Summary}

Graphene \acrshort{gfet} devices provide an excellent transducer for label-free nucleic-acid
detection, but their performance hinges on the capture probe. The biophysical criteria for a
good probe --- conservation, favourable hybridisation thermodynamics and absence of stable
self-structure --- are individually well understood and supported by mature software, yet they
are rarely combined into a single, reproducible and organism-aware pipeline that also ranks the
candidates and validates them structurally. The methodology presented in the next chapter is
designed to fill exactly that gap.
